% !PRINTED 561
\weberpage{561}
%<<CONT
In this, however, the order of summation may be interchanged,
so that we also have

\weq{3}{\eta^{(\lambda)}\eta^{(\mu)} = \sum^{t}\sum^{s} r^{(\lambda + \mu g^{te})g^{se}}}

The sum taken with respect to $s$,

\[
\sum^{s} r^{(\lambda + \mu g^{te})g^{se}}
\]

is, however, itself one of the periods, namely, according to the notation (1), the period

\[
\eta^{(\lambda + \mu g^{te})}.
\]

From (2) and (3) it therefore follows that

\weq{4}{\eta^{(\lambda)}\eta^{(\mu)} = \eta^{(\lambda+\mu)} + \eta^{(\lambda+\mu')} + \eta^{(\lambda+\mu'')} + \cdots}

The right-hand side of this expression contains $f$ terms; among them the same period may of course occur several times;
there may also occur among them the improper period $\eta^{(0)}$,
which is to be set equal to the integer $f$. If one wishes to restore the homogeneous form of the expression, one makes use of
the relation

\[
\eta + \eta_1 + \eta_2 + \cdots + \eta_{e-1} = -1,
\]

which is indeed only another way of writing the equation \S.\,167, (2).

When one has computed according to (4) the expression for a product $\eta\eta_h$,
to which we have already given the form above,

\weq{5}{\eta\eta_h = a_{0,h}\eta + a_{1,h}\eta_1 + \cdots + a_{e-1,h}\eta_{e-1},}

one obtains the product of any two of the periods
by the substitution $(\eta, \eta_k)$,

\weq{6}{\eta_k\eta_{h+k} = a_{0,h}\eta_k + a_{1,h}\eta_{k+1} + \cdots + a_{e-1,h}\eta_{k-1}.}

In order to illustrate these rules by a simple example,
we take $n = 13$. For 13, 2 is a primitive root, and
we can apply the index table given in \S.\,136:

\weq{7}{
\begin{array}{c|cccccccccccc}
I & 0, & 1, & 2, & 3, & 4, & 5, & 6, & 7, & 8, & 9, & 10, & 11\\
\hline
N & 1, & 2, & 4, & 8, & 3, & 6, & 12, & 11, & 9, & 5, & 10, & 7
\end{array}.
}

If we first assume $e = 3$, $f = 4$, we obtain
a cubic resolvent whose roots are

\weq{8}{\begin{aligned}
\eta &= r + r_3 + r_6 + r_9\\
\eta_1 &= r_1 + r_4 + r_7 + r_{10}\\
\eta_2 &= r_2 + r_5 + r_8 + r_{11}
\end{aligned}}

From the table (7) this also gives
%CONT>>
